
\documentclass[border=8pt, multi, tikz]{standalone}
\usepackage{import}
\subimport{../layers/}{init}
\usetikzlibrary{positioning}
\usetikzlibrary{3d}
\usetikzlibrary{calc}
\usepackage[T2A]{fontenc}
\usepackage[utf8]{inputenc}
\usepackage[main=russian,english]{babel}

\def\ConvColor{rgb:yellow,5;red,2.5;white,5}
\def\ConvReluColor{rgb:yellow,5;red,5;white,5}
\def\PoolColor{rgb:red,1;black,0.3}
\def\UnpoolColor{rgb:blue,2;green,1;black,0.3}
\def\FcColor{rgb:blue,5;red,2.5;white,5}
\def\FcReluColor{rgb:blue,5;red,5;white,4}
\def\SoftmaxColor{rgb:magenta,5;black,7}
\def\SumColor{rgb:blue,5;green,15}

\newcommand{\copymidarrow}{\tikz \draw[-Stealth,line width=0.8mm,draw={rgb:blue,4;red,1;green,1;black,3}] (-0.3,0) -- ++(0.3,0);}

\begin{document}
\begin{tikzpicture}
\tikzstyle{connection}=[ultra thick,every node/.style={sloped,allow upside down},draw=\edgecolor,opacity=0.7]
\tikzstyle{copyconnection}=[ultra thick,every node/.style={sloped,allow upside down},draw={rgb:blue,4;red,1;green,1;black,3},opacity=0.7]

%% Входное РЛС-изображение
\node[canvas is zy plane at x=0] (temp) at (-3,0,0)
    {\includegraphics[width=8cm,height=8cm]{../images/sar_example.jpg}};

%% Входной тензор РЛС
\pic[shift={(0,0,0)}] at (0,0,0)
    {Box={
        name=sar_input,
        caption=РЛС,
        xlabel={{1, }},
        zlabel=256,
        fill=\ConvColor,
        height=40,
        width=1,
        depth=40
        }
    };

%% Адаптер РЛС
\pic[shift={(1.5,0,0)}] at (sar_input-east)
    {RightBandedBox={
        name=ccr_sar,
        caption={РЛС + Ли + Собель},
        xlabel={{ 1, 1, 1 }},
        zlabel=256,
        fill=\ConvColor,
        bandfill=\ConvReluColor,
        height=40,
        width={ 2 , 2 , 2 },
        depth=40
        }
    };

\draw [connection] (sar_input-east) -- node {\midarrow} (ccr_sar-west);

%% Хаар-стем
\pic[shift={(1.5,0,0)}] at (ccr_sar-east)
    {Box={
        name=haar_stem,
        caption={Haar-stem\\DWT},
        xlabel={{4, }},
        zlabel=128,
        fill=\ConvColor,
        height=32,
        width=2,
        depth=32
        }
    };

\draw [connection] (ccr_sar-east) -- node {\midarrow} (haar_stem-west);

%% ConvNeXt V2 Этап 1
\pic[shift={(1.5,0,0)}] at (haar_stem-east)
    {RightBandedBox={
        name=enc_s1,
        caption={ConvNeXt\\Этап 1},
        xlabel={{ 96, 96 }},
        zlabel=64,
        fill=\ConvColor,
        bandfill=\ConvReluColor,
        height=25,
        width={ 3 , 3 },
        depth=25
        }
    };

\draw [connection] (haar_stem-east) -- node {\midarrow} (enc_s1-west);

%% ConvNeXt V2 Этап 2
\pic[shift={(1.5,0,0)}] at (enc_s1-east)
    {RightBandedBox={
        name=enc_s2,
        caption={ConvNeXt\\Этап 2},
        xlabel={{ 192, 192 }},
        zlabel=32,
        fill=\ConvColor,
        bandfill=\ConvReluColor,
        height=16,
        width={ 4 , 4 },
        depth=16
        }
    };

\draw [connection] (enc_s1-east) -- node {\midarrow} (enc_s2-west);

%% ConvNeXt V2 Этап 3
\pic[shift={(1.5,0,0)}] at (enc_s2-east)
    {RightBandedBox={
        name=enc_s3,
        caption={ConvNeXt\\Этап 3},
        xlabel={{ 384, 384 }},
        zlabel=16,
        fill=\ConvColor,
        bandfill=\ConvReluColor,
        height=10,
        width={ 5 , 5 },
        depth=10
        }
    };

\draw [connection] (enc_s2-east) -- node {\midarrow} (enc_s3-west);

%% ConvNeXt V2 Этап 4 / Узкое место
\pic[shift={(2,0,0)}] at (enc_s3-east)
    {RightBandedBox={
        name=enc_s4,
        caption={ConvNeXt\\Этап 4\\$8{\times}8$},
        xlabel={{ 768, 768 }},
        zlabel=8,
        fill=\ConvColor,
        bandfill=\ConvReluColor,
        height=6,
        width={ 6 , 6 },
        depth=6
        }
    };

\draw [connection] (enc_s3-east) -- node {\midarrow} (enc_s4-west);

%% PixelShuffle Этап 4  (256 кан., 16x16)
\pic[shift={(2,0,0)}] at (enc_s4-east)
    {Box={
        name=dec_s4,
        caption={PixelShuffle\\Этап 4},
        xlabel={{256, }},
        zlabel=16,
        fill=\UnpoolColor,
        opacity=0.8,
        height=10,
        width=3.5,
        depth=10
        }
    };

\draw [connection] (enc_s4-east) -- node {\midarrow} (dec_s4-west);

%% PixelShuffle Этап 3  (128 кан., 32x32)
\pic[shift={(2,0,0)}] at (dec_s4-east)
    {Box={
        name=dec_s3,
        caption={PixelShuffle\\Этап 3},
        xlabel={{128, }},
        zlabel=32,
        fill=\UnpoolColor,
        opacity=0.8,
        height=16,
        width=3,
        depth=16
        }
    };

\draw [connection] (dec_s4-east) -- node {\midarrow} (dec_s3-west);

%% PixelShuffle Этап 2  (64 кан., 64x64)
\pic[shift={(2,0,0)}] at (dec_s3-east)
    {Box={
        name=dec_s2,
        caption={PixelShuffle\\Этап 2},
        xlabel={{64, }},
        zlabel=64,
        fill=\UnpoolColor,
        opacity=0.8,
        height=25,
        width=2.5,
        depth=25
        }
    };

\draw [connection] (dec_s3-east) -- node {\midarrow} (dec_s2-west);

%% PixelShuffle Этап 1  (32 кан., 128x128)
\pic[shift={(2,0,0)}] at (dec_s2-east)
    {Box={
        name=dec_s1,
        caption={PixelShuffle\\Этап 1},
        xlabel={{32, }},
        zlabel=128,
        fill=\UnpoolColor,
        opacity=0.8,
        height=32,
        width=2,
        depth=32
        }
    };

\draw [connection] (dec_s2-east) -- node {\midarrow} (dec_s1-west);

%% Субполосная голова  (B, 12, 128, 128) = 3 x [LL LH HL HH]
\pic[shift={(1.5,0,0)}] at (dec_s1-east)
    {Box={
        name=subband,
        caption={Субполосы\\LL LH HL HH},
        xlabel={{12, }},
        zlabel=128,
        fill=\SoftmaxColor,
        opacity=0.8,
        height=32,
        width=2,
        depth=32
        }
    };

\draw [connection] (dec_s1-east) -- node {\midarrow} (subband-west);

%% Обратный Хаар (фиксированный, без весов)
\pic[shift={(1.5,0,0)}] at (subband-east)
    {Box={
        name=inv_haar,
        caption={IHaar\\фикс., tanh},
        xlabel={{3, }},
        zlabel=256,
        fill=\ConvColor,
        height=40,
        width=2,
        depth=40
        }
    };

\draw [connection] (subband-east) -- node {\midarrow} (inv_haar-west);

%% Выходное оптическое изображение
\node[canvas is zy plane at x=0] (out_img) at ($(inv_haar-east)+(3,0,0)$)
    {\includegraphics[width=8cm,height=8cm]{../images/opt_example.jpg}};

%% Пропускные связи  (кодировщик -> декодер, синие дуги)
\path (enc_s3-southeast) -- (enc_s3-northeast) coordinate[pos=1.25] (enc_s3-top);
\path (dec_s4-south)     -- (dec_s4-north)     coordinate[pos=1.25] (dec_s4-top);
\draw [copyconnection] (enc_s3-northeast)
    -- node {\copymidarrow} (enc_s3-top)
    -- node {\copymidarrow} (dec_s4-top)
    -- node {\copymidarrow} (dec_s4-north);

\path (enc_s2-southeast) -- (enc_s2-northeast) coordinate[pos=1.25] (enc_s2-top);
\path (dec_s3-south)     -- (dec_s3-north)     coordinate[pos=1.25] (dec_s3-top);
\draw [copyconnection] (enc_s2-northeast)
    -- node {\copymidarrow} (enc_s2-top)
    -- node {\copymidarrow} (dec_s3-top)
    -- node {\copymidarrow} (dec_s3-north);

\path (enc_s1-southeast) -- (enc_s1-northeast) coordinate[pos=1.25] (enc_s1-top);
\path (dec_s2-south)     -- (dec_s2-north)     coordinate[pos=1.25] (dec_s2-top);
\draw [copyconnection] (enc_s1-northeast)
    -- node {\copymidarrow} (enc_s1-top)
    -- node {\copymidarrow} (dec_s2-top)
    -- node {\copymidarrow} (dec_s2-north);

\end{tikzpicture}
\end{document}
